\documentclass[11pt]{article}
\usepackage[margin=.7in,top=.55in,bottom=.55in]{geometry}
\usepackage{amsmath,amssymb,amsthm,hyperref}
\renewcommand{\normalsize}{\fontsize{10}{12}\selectfont}
\DeclareMathSizes{10}{10}{8}{6}
\hypersetup{colorlinks=true,linkcolor=blue,urlcolor=blue,citecolor=blue}
\setlength{\emergencystretch}{2em}
\newtheorem{theorem}{Theorem}
\newtheorem{lemma}{Lemma}
\newcommand{\tr}{\operatorname{tr}}
\newcommand{\Sk}[2]{S_{#1}(#2)}
\begin{document}
\begin{center}
{\Large\bfseries Strict connection--Dirac spectral dominance\\for every finite graph}\\[.5em]
Research proof and formalization draft\\
October 26, 2026
\end{center}
\begin{abstract}
Let $L$ be the simplex-intersection connection matrix of a nonempty finite
graph, viewed as its vertex-and-edge simplicial complex, and let $D$ be its
signed simplicial Dirac matrix. For every $1\le k\le v+m$ we prove
$\Sk{k}{L}-\Sk{k}{D}\ge\sqrt2-1>0$, where $S_k$ sums the $k$ largest
eigenvalues with multiplicity, $v$ is the number of vertices, and $m$ the
number of edges. Isolated vertices and disconnected graphs are allowed.
This establishes, with a uniform strict gap, the dimension-at-most-one case
of Knill's connection--Dirac conjecture. The higher-dimensional conjecture
remains outside the result. The proof uses incidence matrices, orthonormal
eigenvector lifts, the Ky Fan variational principle, and scalar inequalities.
\end{abstract}

\section{Statement and source matrices}
For a real symmetric $N\times N$ matrix $A$ write
$\lambda_1(A)\ge\cdots\ge\lambda_N(A)$ and
$\Sk{k}{A}=\sum_{i=1}^k\lambda_i(A)$.
Knill \cite[Conjecture~1, \S5.3]{K26} asks whether the connection matrix
weakly dominates the Dirac matrix in this sense for every finite abstract
simplicial complex. We prove the full graph case.

Let $G$ be a finite simple graph with $v\ge1$ vertices and $m$ edges.
Its nonempty simplices are precisely its vertices and edges; no larger
clique is filled. Thus a triangle means its boundary complex.
For simplices $s,t$ put $L_{st}=1$ if $s\cap t\ne\varnothing$, and $0$
otherwise. Let $d$ be an oriented simplicial boundary matrix and
$D=d+d^{\mathsf T}$.
Changing edge orientations conjugates $D$ by a diagonal orthogonal matrix.
\begin{theorem}\label{main}
For every such graph and every $1\le k\le v+m$,
\[
 \Sk{k}{L}-\Sk{k}{D}\ge\sqrt2-1>0.
\]
The constant is a uniform lower bound; no optimality is asserted.
\end{theorem}

Order vertices before edges. Let $U$ be unsigned vertex--edge incidence,
and $C$ signed incidence, with one $+1$ and one $-1$ in each column.
Write $\delta_i$ for the vertex degrees, $\Delta=\operatorname{diag}(\delta_i)$,
$Q=UU^{\mathsf T}$, and $H=CC^{\mathsf T}$. The literal matrices are
\begin{equation}\label{blocks}
 L=\begin{pmatrix}I&U\\U^{\mathsf T}&U^{\mathsf T}U-I\end{pmatrix},
 \qquad
 D=\begin{pmatrix}0&C\\C^{\mathsf T}&0\end{pmatrix}.
\end{equation}
In the edge block, distinct simple edges share at most one vertex and
$(U^{\mathsf T}U)_{ee}=2$. Consequently \eqref{blocks} is exactly the
intersection definition. Also
\begin{equation}\label{degrees}
 Q_{ii}=\delta_i,\qquad \tr Q=\tr H=2m,\qquad H\preceq2\Delta,
\end{equation}
the last inequality following from
$(x_i-x_j)^2\le2x_i^2+2x_j^2$ on every edge.
Here $\preceq$ denotes ordinary positive-semidefinite order.
Let $T_k$ be the sum of the $k$ largest degrees, for $k\le v$.

\section{The connection eigenvector lift}
Set
\[
 \phi(q)=\frac{q+\sqrt{q^2+4}}2,
 \qquad e(q)=\frac{\sqrt{q^2+4}}2-1.
\]
For $q\ge0$ one has $\phi(q)>0$,
$\phi(q)^2-q\phi(q)=1$, and
$\phi(q)=1+q/2+e(q)$ with $e(q)\ge0$.
If $Qu=qu$, direct substitution into \eqref{blocks} shows that
\begin{equation}\label{lift}
 \bigl((\phi(q)+1)u,\ \phi(q)U^{\mathsf T}u\bigr)
\end{equation}
is an $L$-eigenvector with eigenvalue $\phi(q)$.
This formula also works at $q=0$. It is nonzero for $u\ne0$.
An orthonormal $Q$-eigenbasis yields mutually orthogonal vectors
\eqref{lift}, because
$\langle U^{\mathsf T}u_i,U^{\mathsf T}u_j\rangle
=\langle u_i,Qu_j\rangle=q_j\delta_{ij}$.
Normalize these vectors and apply the Ky Fan principle. If
$q_1\ge\cdots\ge q_v\ge0$ are the eigenvalues of $Q$, then
\begin{equation}\label{positive-lower}
 \Sk{k}{L}\ge\sum_{i=1}^k\phi(q_i)\quad(1\le k\le v).
\end{equation}
The incidence reduction is consistent with Knill's earlier hydrogen
identity \cite[Theorem~1]{K18}; the block reduction itself is a known
ingredient, not an additional novelty claim.

The diagonal Ky Fan inequality gives
$\sum_{i=1}^kq_i\ge T_k$.
If $m>0$, choose an edge $\{a,b\}$ and $z=e_a+e_b$. Then
$\|z\|^2=2$ and $z^{\mathsf T}Qz=\delta_a+\delta_b+2\ge4$,
so $q_1\ge2$. Since $e$ increases on $[0,\infty)$,
\begin{equation}\label{low-L}
 \Sk{k}{L}\ge k+T_k/2+e(q_1)
 \ge k+T_k/2+\sqrt2-1\qquad(1\le k\le v).
\end{equation}

\section{Upper bounds for the signed Dirac matrix}
Consider any orthonormal family of positive $D$-eigenvectors
$w_i=(a_i,b_i)$ with eigenvalues $\mu_i>0$.
The block equations are
$Cb_i=\mu_i a_i$ and $C^{\mathsf T}a_i=\mu_i b_i$.
Adjointness gives
\[
 \mu_i\langle b_i,b_j\rangle=\mu_j\langle a_i,a_j\rangle,
 \qquad
 \mu_j\langle b_i,b_j\rangle=\mu_i\langle a_i,a_j\rangle.
\]
Adding and using $\mu_i+\mu_j>0$ shows that the two inner products
are equal. Their sum is $\delta_{ij}$, hence
$u_i=\sqrt2a_i$ is an orthonormal family in $\mathbb R^v$ and
$Hu_i=\mu_i^2u_i$. This treats repeated eigenvalues directly.

Take the positive eigenvalues among the first $k$ eigenvalues of $D$,
and denote their number by $r$. Dropping the nonpositive entries increases
the sum. For $k\le v$ we have $r\le k$; using \eqref{degrees}, Ky Fan,
and nonnegativity of the degrees gives
\[
 \sum_{i=1}^r\mu_i^2
 =\sum_{i=1}^r\langle u_i,Hu_i\rangle
 \le2T_r\le2T_k.
\]
Cauchy--Schwarz therefore proves
\begin{equation}\label{low-D}
 \Sk{k}{D}\le\sum_{i=1}^r\mu_i
 \le\sqrt{r\sum_{i=1}^r\mu_i^2}\le\sqrt{2kT_k}.
\end{equation}
For arbitrary $k\le v+m$, orthonormality instead gives $r\le v$ and
$\sum_{i=1}^r\mu_i^2\le\tr H=2m$. The scalar inequality
$t\le(t^2+1)/2$ for $t\ge0$ yields
\begin{equation}\label{global-D}
 \Sk{k}{D}\le\frac{\sum_{i=1}^r\mu_i^2+r}{2}\le m+v/2.
\end{equation}
This argument only needs upper bounds; it makes no assumption about
the zero multiplicity in the full Dirac spectrum.

\section{Completion of the comparison}
Suppose first $m>0$ and $k\le v$. Since
\[
 (k+T_k/2)^2-2kT_k=(k-T_k/2)^2\ge0,
\]
we have $k+T_k/2\ge\sqrt{2kT_k}$. Subtract
\eqref{low-D} from \eqref{low-L} to obtain the claimed gap.

For the remaining indices we use an alternative to a full inertia calculation.
The quadratic identity, for $x=(a,b)$,
\[
 x^{\mathsf T}(L+I)x=\|a\|^2+\|a+Ub\|^2\ge0
\]
proves $\lambda_i(L)\ge-1$.
Put $g=\sum_{i=1}^v\phi(q_i)^{-1}>0$.
The identity $\phi(q)-\phi(q)^{-1}=q$ and \eqref{degrees} give
$\sum_i\phi(q_i)=2m+g$. By Cauchy--Schwarz,
\[
 v^2\le\Bigl(\sum_i\phi(q_i)\Bigr)
             \Bigl(\sum_i\phi(q_i)^{-1}\Bigr)=(2m+g)g.
\]
Simplicity gives $2m\le v(v-1)$, and $v\ge1$ implies $2m+1\le v^2$.
If $g<1$, then $(2m+g)g<2m+1$, a contradiction. Thus $g\ge1$ and
\eqref{positive-lower} gives $\Sk{v}{L}\ge2m+1$.
For $v\le k\le v+m$ it follows that
\[
 \Sk{k}{L}\ge2m+v+1-k.
\]
On the other hand, the diagonal of $L$ consists of ones, so Ky Fan gives
$\Sk{k}{L}\ge k$. Averaging these two lower bounds yields
\[
 \Sk{k}{L}\ge m+v/2+1/2.
\]
Together with \eqref{global-D}, this proves a gap at least
$1/2>\sqrt2-1$ for these indices.
Finally, if $m=0$, then $L=I_v$ and $D=0$, so the gap is $k\ge1$.
All cases of Theorem~\ref{main} follow.

\section{Formal development and scope}
\textbf{Development status.} A fresh isolated replay of ten accompanying
Lean modules verifies 135 declarations, including the incidence identities,
the variational principle for signed spectral arrays, the normalized
connection eigenvector family, both connection lower bounds, the edgeless
comparison, and the bijection with the downward-closed graph face family.
The Dirac bounds and combined theorem endpoint are still being formalized.
Thus this draft does not yet claim a complete formal certificate.
The replay uses Lean~4.34.1 and mathlib revision \texttt{d13f23b7}, treats
warnings as errors, and checks every declaration's axioms. Only
\texttt{propext}, \texttt{Classical.choice}, and \texttt{Quot.sound} occur.

The contribution is one general comparison for every graph, rather than
separate claims for its helpers, special cases, or constant.
The source problem for complexes with higher-dimensional faces is untouched.
Knill's hydrogen identity, incidence formulas, spectral variational facts,
and scalar inequalities are credited as established ingredients.
The packet records the primary sources and a bounded prior-art investigation;
no worldwide priority claim or complete literature census is asserted.

\begin{thebibliography}{9}
\bibitem{K26} O.~Knill,
\emph{Remarks about connection and Dirac matrices},
arXiv:2601.18071v2 (29 July 2026),
\url{https://arxiv.org/abs/2601.18071}.
\bibitem{K18} O.~Knill,
\emph{The hydrogen identity for Laplacians},
arXiv:1803.01464v1 (2018),
\url{https://arxiv.org/abs/1803.01464}.
\end{thebibliography}
\end{document}
